% ===== Experiment 2: filled from figures/macros2.tex =====
\subsection{Why intent-aligned allocation is locally rational --- and when it stops being so}
\label{sec:theory}

Within a single decision, Experiment 1's allocation is defensible. Let $A$ be
the intended action and $m_A$ the memories backing it. Verifying $m \in m_A$
has immediate value: if $m$ is wrong, the agent changes what it is about to
do. Verifying a belief behind an action the agent is \emph{not} taking has,
within the episode, an expected value of zero --- the belief is never
exercised. A myopic agent that spends its budget on $m_A$ is optimizing
correctly under the objective ``make today's decision well.''

The value of verifying memory $m$ is better written
\[
\mathrm{EVOI}(m) \;=\; \textstyle\sum_{c} P(c)\, P(\text{$m$ load-bearing in } c)\,
P(\text{$m$ corrupted})\, L(m, c),
\]
where $c$ ranges over the \emph{future contexts the belief will survive into},
and $L$ is the loss averted by catching the corruption before acting in $c$.
A single-decision agent truncates the sum to the present context; a
long-horizon agent should not, because durable memory outlives the decision
that happened to be on the desk when the budget was spent. The experiments so
far measure the truncated policy. Experiment 2 measures its cost: it makes one
of the discounted future contexts arrive.

\subsection{Experiment 2: the unexamined belief becomes load-bearing}

Phase 1 reproduces Experiment 1 exactly (budget $k{=}2$, target the pricing
memory). Then, in the same session, the situation shifts: the competitor cuts
prices again, churn among price-sensitive SMB accounts doubles, and leadership
wants a response \emph{today} --- with the verification budget spent. The
corrupted belief is now squarely decision-relevant, and whether its lost
caveat is in the agent's context depends entirely on how the phase-1 budget
was allocated. We run \PTwoMainN{} two-phase episodes (six models $\times$
three arms $\times$ 25), pre-registered as H6--H8 with scoring fixed in
advance. \emph{Reliance} on the corrupted lesson means: promotional pricing
chosen, the corrupted memory cited as evidence, and no uncertainty preserved.
The strict \emph{harmful} rule additionally requires more than a small guarded
test and that the corrupted memory was never verified.

\paragraph{The caveat's visible presence prevents what its silent absence
permits (H7).}
Under drift, phase-2 reliance occurs in \PTwoRelDriftK{}/\PTwoRelDriftN{}
episodes and the strict harmful rule fires in
\PTwoHarmDriftK{}/\PTwoHarmDriftN{} --- the first behavioral failures anywhere
in this study. In the clean arm, with the caveat visible in the inherited
memory itself, reliance is \PTwoRelCleanK{}/\PTwoRelCleanN{} and harm is
\PTwoHarmCleanK{}/\PTwoHarmCleanN{} (clean vs.\ drifted harm,
$p = \PTwoHarmFisherP$). Nothing about the phase-2 pressure forces the
mistake: models facing the same ultimatum with the caveat intact hedge, guard,
or choose differently.

\paragraph{Phase-1 verification protects phase 2 (H6).}
The randomized comparison is between the drifted and drifted-triage arms:
identical lineage, identical phase-2 ultimatum, differing only in the
one-sentence triage invariant that (Experiment 1) lifts verification of the
corrupted memory. Reliance moves from \PTwoRelDriftK{}/\PTwoRelDriftN{} to
\PTwoRelTriageK{}/\PTwoRelTriageN{} (CMH stratified by model,
$p = \PTwoArmCMHP$), and harm from \PTwoHarmDriftK{} to \PTwoHarmTriageK{}
episodes. The self-selected conditional points the same way --- reliance in
\PTwoCondUK{}/\PTwoCondUN{} of drifted episodes whose budget missed the
corrupted memory, against \PTwoCondVK{}/\PTwoCondVN{} of those that reached it
(CMH $p = \PTwoCondCMHP$) --- but models that verify differ from models that
do not, so we lead with the randomized contrast.

\paragraph{Knowing is not allocating (H8).}
A separate probe hands each model the drifted lineage and asks it to rate how
well-supported each inherited memory is for the stated objective --- no
budget, no decision, order randomized as always. If low verification reflected
a failure to \emph{notice} the corrupted memory's weakness, the ratings should
miss it too. Table~\ref{tab:probe} reports, per model, how often the corrupted
memory is named least-supported in the audit against how often the same
model's phase-1 budget actually reached it.

\begin{table}[t]
\centering
\small
\begin{tabular}{lccc}
\toprule
Model & Named least-supported & Strictly lowest rating & Verified at $k{=}2$ (Exp.~1) \\
\midrule
Claude Opus 5    & \PRBopusNamedK{}/\PRBopusN{}     & \PRBopusStrictK{}/\PRBopusN{}   & \VopusDriftK{}/\VopusDriftN{} \\
Claude Sonnet 5  & \PRBsonnetNamedK{}/\PRBsonnetN{} & \PRBsonnetStrictK{}/\PRBsonnetN{} & \VsonnetDriftK{}/\VsonnetDriftN{} \\
Claude Haiku 4.5 & \PRBhaikuNamedK{}/\PRBhaikuN{}   & \PRBhaikuStrictK{}/\PRBhaikuN{} & \VhaikuDriftK{}/\VhaikuDriftN{} \\
GPT-5.6 Sol      & \PRBsolNamedK{}/\PRBsolN{}       & \PRBsolStrictK{}/\PRBsolN{}     & \VsolDriftK{}/\VsolDriftN{} \\
GPT-5.6 Terra    & \PRBterraNamedK{}/\PRBterraN{}   & \PRBterraStrictK{}/\PRBterraN{} & \VterraDriftK{}/\VterraDriftN{} \\
GPT-5.6 Luna     & \PRBlunaNamedK{}/\PRBlunaN{}     & \PRBlunaStrictK{}/\PRBlunaN{}   & \VlunaDriftK{}/\VlunaDriftN{} \\
\bottomrule
\end{tabular}
\caption{The knowledge--allocation gap. Asked to audit the drifted lineage
with no budget at stake, models rate the corrupted memory least-supported far
more often than their own scarce verification reaches it.}
\label{tab:probe}
\end{table}
